\section{Determinants of small matrices}

We begin with explicit formulas. These formulas are the definitions for the small sizes considered in this section.

\subsection{The determinant of a \(1\times1\) matrix}

For a matrix containing one entry, we define
\[
\det\begin{pmatrix}a\end{pmatrix}=a.
\]
The determinant is simply the entry itself.

This convention will later be needed when larger determinants are reduced to smaller ones.

\subsection{The determinant of a \(2\times2\) matrix}

\begin{definitionbox}
For
\[
A=\begin{pmatrix}a&b\\c&d\end{pmatrix},
\]
the determinant of \(A\) is
\[
\det A=ad-bc.
\]
We also write
\[
\begin{vmatrix}a&b\\c&d\end{vmatrix}=ad-bc.
\]
\end{definitionbox}

The products \(ad\) and \(bc\) use the two diagonals of the matrix. The order matters: the second product is subtracted from the first.

\begin{examplebox}
Let
\[
A=\begin{pmatrix}3&1\\2&5\end{pmatrix}.
\]
Then
\[
\det A=3\cdot5-1\cdot2=15-2=13.
\]
\end{examplebox}

\begin{examplebox}
Let
\[
B=\begin{pmatrix}1&2\\3&6\end{pmatrix}.
\]
Then
\[
\det B=1\cdot6-2\cdot3=6-6=0.
\]
For the moment, this is simply the result of applying the definition. In the next section we will examine why proportional rows force this cancellation.
\end{examplebox}

A determinant is sensitive to the positions of the entries. For example,
\[
\det\begin{pmatrix}1&2\\3&4\end{pmatrix}=4-6=-2,
\]
whereas
\[
\det\begin{pmatrix}3&4\\1&2\end{pmatrix}=6-4=2.
\]
The same four numbers appear, but the rows have been exchanged and the sign has changed.

\subsection{The determinant of a \(3\times3\) matrix}

For a \(3\times3\) matrix we again begin with an explicit formula.

\begin{definitionbox}
For
\[
A=\begin{pmatrix}
a&b&c\\
d&e&f\\
g&h&i
\end{pmatrix},
\]
we define
\[
\det A=aei+bfg+cdh-ceg-bdi-afh.
\]
\end{definitionbox}

One way to remember this formula is Sarrus' rule. Repeat the first two columns:
\[
\begin{array}{ccc|cc}
a&b&c&a&b\\
d&e&f&d&e\\
g&h&i&g&h
\end{array}
\]
Add the products along the three downward diagonals and subtract the products along the three upward diagonals:
\[
\det A=(aei+bfg+cdh)-(ceg+bdi+afh).
\]

\begin{remarkbox}
Sarrus' rule is a mnemonic for the \(3\times3\) formula. It is not a rule for determinants of size \(4\times4\) or larger.
\end{remarkbox}

\begin{examplebox}
Let
\[
A=\begin{pmatrix}
1&2&0\\
-1&3&4\\
2&0&5
\end{pmatrix}.
\]
Then
\[
\begin{aligned}
\det A
&=(1\cdot3\cdot5+2\cdot4\cdot2+0\cdot(-1)\cdot0)\\
&\quad -(0\cdot3\cdot2+2\cdot(-1)\cdot5+1\cdot4\cdot0)\\
&=(15+16+0)-(0-10+0)\\
&=41.
\end{aligned}
\]
\end{examplebox}

\begin{examplebox}
Consider
\[
C=\begin{pmatrix}
1&2&0\\
2&4&0\\
3&1&5
\end{pmatrix}.
\]
Using the formula,
\[
\det C
=1\cdot4\cdot5+2\cdot0\cdot3+0\cdot2\cdot1
-0\cdot4\cdot3-2\cdot2\cdot5-1\cdot0\cdot1=0.
\]
The first two rows are proportional. The cancellation is not an accident; it is one of the basic properties of determinants.
\end{examplebox}

\subsection{What should be learned from the formulas?}

At this stage there are two separate skills.
\begin{itemize}
\item Apply the correct formula without losing signs or terms.
\item Compare examples and look for regular behaviour.
\end{itemize}

The second skill leads to the next section. Instead of computing many unrelated examples, we will use the formulas to establish rules that hold for whole families of matrices.
